\documentclass[11pt,a4paper]{article}
\usepackage[margin=25mm]{geometry}
\usepackage{booktabs}
\usepackage{amsmath}
\usepackage{graphicx}
\usepackage{hyperref}
\usepackage{microtype}
\usepackage{parskip}
\usepackage{caption}
\usepackage{multirow}
\usepackage{array}
\usepackage{float}
\usepackage{enumitem}

\hypersetup{colorlinks=true,linkcolor=black,citecolor=black,urlcolor=black}

\title{\textbf{Project Assessment 5}\\
Cross-Domain Replicated Analysis of Selection Perturbative\\
Hyper-Heuristics with Explainable AI}
\author{Dreas Vermaak u22497618, Liam van Kasterop u22539761,\\
Katleho Moeletsi 19061430, Emmanuel Ramphabana u25721373\\
Department of Computer Science}
\date{October 2026}

\begin{document}
\maketitle

\section{Introduction}

This assessment addresses four feedback items from Assessment~4: insufficient
experimental replication, an incorrect claim of comparability with published
CHeSC results, inconsistent cross-technique SHAP magnitude comparisons, and
the absence of explicit research questions to guide the analysis. The
hyper-heuristic logic is unchanged; the improvements are in experimental
methodology, statistical rigour, and explanation validity.

\subsection{Research Questions}

\begin{itemize}
  \item \textbf{RQ1}: Within each domain, does any technique outperform the Simple Random baseline (SP-SR-LA) with statistical support?
  \item \textbf{RQ2}: Which technique performs best across domains?
  \item \textbf{RQ3}: Which state features drive acceptance within each technique, and are these explanations stable across seeds?
  \item \textbf{RQ4}: Do the features that drive acceptance differ between selection strategies, acceptance strategies, and search topologies?
\end{itemize}

\subsection{Techniques}

\begin{table}[H]
\centering
\begin{tabular}{llll}
\toprule
Technique & Topology & Selection & Acceptance \\
\midrule
SP-SR-LA & Single-point & Simple Random & Late Acceptance, $L{=}50$ \\
SP-CF-LA & Single-point & Choice Function & Late Acceptance, $L{=}50$ \\
SP-CF-AILTA & Single-point & Choice Function & AILTA \\
MP-CF & Multi-point (pop=10) & Choice Function & Population replacement \\
\bottomrule
\end{tabular}
\end{table}

\textbf{Selection strategies.}
\emph{Simple Random} selects a low-level heuristic uniformly at random each
iteration, serving as a non-adaptive baseline.
The \emph{Modified Choice Function} (Drake, \"{O}zcan \& Burke, 2012) scores each
heuristic~$h$ as
$\mathrm{CF}(h) = \varphi\,f_1[h] + \varphi\,f_2[\textit{prev}][h]
               + (1{-}\varphi)\,f_3[h]$,
where $f_1$ accumulates the improvement per unit CPU time of~$h$ applied alone
(exponentially decayed), $f_2$ captures pair-sequence synergy between the
previous heuristic and~$h$, and $f_3$ is the elapsed time since~$h$ was last
applied (favouring idle heuristics). The adaptive weight~$\varphi$ snaps to
$0.99$ on any improving move and decays by $0.01$ per non-improving step
(floor $0.01$), shifting the balance from exploitation to exploration as
the search stagnates. The first~$n$ iterations (one per heuristic) apply each
heuristic once in round-robin to initialise $f_1$.

\textbf{Acceptance strategies.}
\emph{Late Acceptance} (Burke \& Bykov 2017) maintains a circular buffer of
length $L{=}50$ storing past incumbent objectives. A candidate is accepted if
it improves the current incumbent \emph{or} is no worse than the entry
$L$~steps ago; the buffer is tightened monotonically (never relaxed). All
slots are initialised to $+\infty$, so the first $L$~iterations accept
everything.
\emph{AILTA} is a simplified linear-decay variant of the Adaptive
Iteration-Limited Threshold Accepting strategy used in AdapHH
(M{\i}s{\i}r \& \"{O}zcan 2014). The full AILTA reheats the threshold after
prolonged stagnation; our variant omits reheating and decays the threshold
linearly: $\theta(t) = 0.02\,|f_0|\,(1 - t/T)$, where $f_0$ is the starting
objective. A candidate is accepted if its objective exceeds the incumbent by
at most~$\theta(t)$. At the final iteration, $\theta = 0$ and the strategy
degenerates to improving-or-equal acceptance. The threshold is scaled to each
domain's starting objective, avoiding the need to tune an absolute constant
per domain.

\textbf{Multi-point topology.}
MP-CF maintains a population of $p{=}10$ solutions. Each iteration, a binary
tournament selects the \emph{worst} individual as the replacement target, and
a second tournament (excluding the target) selects the \emph{best} as a
potential crossover partner. A candidate replaces the target if and only if it
is strictly better \emph{and} the target is not the population-best (elitism
guard). There is no separate acceptance module; the replacement rule serves
as the acceptance criterion. A single shared Choice Function instance scores
heuristics across all individuals.

\textbf{Low-level heuristic sets.}
Each HyFlex domain exposes a fixed set of heuristics discovered at runtime
via \texttt{getNumberOfHeuristics()}: SAT (11), BinPacking (8), FlowShop (15),
TSP (13), VRP (10), PersonnelScheduling (12). These include mutation,
ruin-and-recreate, local search, and crossover operators; the hyper-heuristic
treats them as black boxes identified only by index and HyFlex type class.

\section{Experimental Setup}

\subsection{Real HyFlex Domains}

All experiments use the real CHeSC HyFlex jar files placed in
\texttt{lib/hyflex/}. All 240 run manifests report \texttt{stub:~false}.
The six domains are SAT, BinPacking, FlowShop, TSP, VRP, and
PersonnelScheduling.

\textbf{Limitation}: These results are not directly comparable to published CHeSC
rankings. The CHeSC competition used a normalised 10-minute wall-clock time limit
(scaled by a hardware benchmark), competition instances, and 31 independent runs per
technique per domain. This study uses a per-domain iteration budget, instance~0 only,
and 10 seeds. The real HyFlex domains are the same, but the experimental protocol
differs.

\subsection{Budget Policy}

The budget is defined as a count of heuristic applications (calls to
\texttt{applyHeuristic()}). All four techniques within a domain receive the same
budget. MP-CF with population~10 receives the same total number of applications as
the single-point techniques, not 10 times as many.

Per-domain budgets are set from convergence curves on a pilot run (Assessment~4,
seed~42, 100{,}000 iterations for non-PS domains, 10{,}000 for PersonnelScheduling).
For each domain, the budget is the smallest tier in
\{5{,}000; 10{,}000; 20{,}000; 50{,}000; 100{,}000\} that exceeds the iteration at
which all four techniques reach 95\% of their total improvement.

\begin{table}[H]
\centering
\begin{tabular}{lrl}
\toprule
Domain & Budget & 95\% convergence iteration \\
\midrule
SAT & 50{,}000 & 28{,}536 \\
BinPacking & 100{,}000 & 94{,}670 \\
FlowShop & 100{,}000 & 91{,}405 \\
TSP & 10{,}000 & 9{,}834 \\
VRP & 50{,}000 & 29{,}002 \\
PersonnelScheduling & 5{,}000 & 1{,}212 \\
\bottomrule
\end{tabular}
\end{table}

PersonnelScheduling's budget is naturally low because it converges rapidly. Profiling
(Task~0) showed that 94.8\% of PersonnelScheduling wall time is spent in HyFlex
\texttt{LOCAL\_SEARCH} calls (mean 484\,ms per application), with instrumentation
contributing only 1.4\%. No optimisation of our code can reduce this cost.

\subsection{Replication}

Each cell (technique $\times$ domain) is run with 10 independent seeds: 42, 123, 456,
789, 1024, 2048, 3141, 4242, 5555, 9999. Each run uses a fresh problem object, fresh
selector state, and a fresh RNG seeded with the run's seed. Total: 240 runs, all
completing successfully. Total wall time: 20 hours 14 minutes.

\subsection{Improvement over Assessment 4}

\begin{table}[H]
\centering
\begin{tabular}{lll}
\toprule
Aspect & Assessment 4 & Assessment 5 \\
\midrule
Seeds per cell & 1 & 10 \\
Total runs & 24 & 240 \\
Statistical tests & None & Wilcoxon, Friedman, A12 \\
Budget & Fixed 100{,}000 all domains & Per-domain from convergence \\
SHAP surrogates & Pooled across domains & Per-domain, per-seed \\
SHAP comparison & Raw magnitudes & Ranks and shares \\
CHeSC claim & Claimed comparability & Explicitly disclaimed \\
\texttt{feat\_llh\_count} & Included (domain label) & Excluded \\
\bottomrule
\end{tabular}
\end{table}

\section{Results: Within-Domain Performance (RQ1)}

\subsection{Summary Statistics}

\textbf{SAT} (budget: 50{,}000)

\begin{table}[H]
\centering
\small
\begin{tabular}{lrrrrrr}
\toprule
Technique & Median & IQR & Mean & SD & Min & Max \\
\midrule
SP-SR-LA & 12.50 & [9.50, 15.75] & 12.50 & 3.78 & 6.0 & 19.0 \\
SP-CF-LA & 11.50 & [8.50, 13.75] & 12.20 & 4.18 & 7.0 & 20.0 \\
SP-CF-AILTA & \textbf{4.00} & [4.00, 4.75] & 4.30 & 1.16 & 3.0 & 7.0 \\
MP-CF & 44.00 & [42.25, 49.00] & 43.80 & 8.31 & 30.0 & 56.0 \\
\bottomrule
\end{tabular}
\end{table}

\textbf{BinPacking} (budget: 100{,}000)

\begin{table}[H]
\centering
\small
\begin{tabular}{lrrrr}
\toprule
Technique & Median & IQR & Mean & SD \\
\midrule
SP-SR-LA & 0.01140 & [0.0110, 0.0118] & 0.0113 & 0.0006 \\
SP-CF-LA & 0.01190 & [0.0114, 0.0124] & 0.0119 & 0.0007 \\
SP-CF-AILTA & 0.02080 & [0.0191, 0.0215] & 0.0205 & 0.0020 \\
MP-CF & \textbf{0.01140} & [0.0112, 0.0115] & 0.0114 & 0.0002 \\
\bottomrule
\end{tabular}
\end{table}

Note: SP-SR-LA and MP-CF are tied at median 0.0114. MP-CF is not significantly
better than SP-SR-LA ($p = 0.557$; see \S3.2), so this ranking is not meaningful.

\textbf{FlowShop} (budget: 100{,}000)

\begin{table}[H]
\centering
\small
\begin{tabular}{lrrrr}
\toprule
Technique & Median & IQR & Mean & SD \\
\midrule
SP-SR-LA & 6353.0 & [6347, 6366] & 6354.3 & 16.6 \\
SP-CF-LA & 6364.0 & [6346, 6374] & 6362.1 & 30.0 \\
SP-CF-AILTA & \textbf{6313.5} & [6303, 6325] & 6311.4 & 18.5 \\
MP-CF & 6369.0 & [6364, 6390] & 6375.5 & 18.3 \\
\bottomrule
\end{tabular}
\end{table}

\textbf{TSP} (budget: 10{,}000)

\begin{table}[H]
\centering
\small
\begin{tabular}{lrrrr}
\toprule
Technique & Median & IQR & Mean & SD \\
\midrule
SP-SR-LA & 51{,}039 & [50{,}607, 52{,}188] & 51{,}206 & 1{,}070 \\
SP-CF-LA & 51{,}558 & [51{,}198, 52{,}633] & 51{,}640 & 1{,}257 \\
SP-CF-AILTA & 51{,}204 & [50{,}231, 51{,}567] & 50{,}996 & 1{,}036 \\
MP-CF & \textbf{49{,}826} & [49{,}743, 50{,}100] & 50{,}010 & 634 \\
\bottomrule
\end{tabular}
\end{table}

\textbf{VRP} (budget: 50{,}000)

\begin{table}[H]
\centering
\small
\begin{tabular}{lrrrr}
\toprule
Technique & Median & IQR & Mean & SD \\
\midrule
SP-SR-LA & 6{,}133 & [5{,}896, 6{,}218] & 6{,}123 & 669 \\
SP-CF-LA & 7{,}132 & [7{,}020, 7{,}164] & 7{,}031 & 584 \\
SP-CF-AILTA & 5{,}215 & [5{,}188, 5{,}234] & 5{,}299 & 289 \\
MP-CF & \textbf{5{,}161} & [5{,}148, 5{,}197] & 5{,}177 & 47 \\
\bottomrule
\end{tabular}
\end{table}

\textbf{PersonnelScheduling} (budget: 5{,}000)

\begin{table}[H]
\centering
\small
\begin{tabular}{lrrrr}
\toprule
Technique & Median & IQR & Mean & SD \\
\midrule
SP-SR-LA & 3{,}334 & [3{,}315, 3{,}340] & 3{,}330 & 24.3 \\
SP-CF-LA & 3{,}336 & [3{,}323, 3{,}355] & 3{,}340 & 24.8 \\
SP-CF-AILTA & \textbf{3{,}301} & [3{,}291, 3{,}311] & 3{,}303 & 14.2 \\
MP-CF & 3{,}341 & [3{,}336, 3{,}345] & 3{,}340 & 7.0 \\
\bottomrule
\end{tabular}
\end{table}

\subsection{Wilcoxon Signed-Rank Tests vs SP-SR-LA Baseline}

All tests are two-sided, paired by seed, with Holm correction within each domain.
$A_{12} < 0.5$ means the tested technique is better (lower objective) than SP-SR-LA.

\begin{table}[H]
\centering
\small
\begin{tabular}{llrrl}
\toprule
Domain & Technique & $p$ (corrected) & $A_{12}$ & Significant? \\
\midrule
SAT & SP-CF-LA & 0.992 & 0.455 & No \\
SAT & SP-CF-AILTA & \textbf{0.006} & \textbf{0.01} & \textbf{Yes} \\
SAT & MP-CF & \textbf{0.006} & 1.00 & \textbf{Yes} (worse) \\
\midrule
BinPacking & SP-CF-LA & \textbf{0.006} & 0.90 & \textbf{Yes} (worse) \\
BinPacking & SP-CF-AILTA & \textbf{0.006} & 1.00 & \textbf{Yes} (worse) \\
BinPacking & MP-CF & 0.557 & 0.45 & No \\
\midrule
FlowShop & SP-CF-LA & 0.541 & 0.59 & No \\
FlowShop & SP-CF-AILTA & \textbf{0.012} & \textbf{0.03} & \textbf{Yes} \\
FlowShop & MP-CF & 0.059 & 0.815 & No \\
\midrule
TSP & SP-CF-LA & 1.000 & 0.61 & No \\
TSP & SP-CF-AILTA & 1.000 & 0.45 & No \\
TSP & MP-CF & 0.059 & 0.18 & No \\
\midrule
VRP & SP-CF-LA & \textbf{0.020} & 0.81 & \textbf{Yes} (worse) \\
VRP & SP-CF-AILTA & \textbf{0.020} & \textbf{0.19} & \textbf{Yes} \\
VRP & MP-CF & \textbf{0.018} & \textbf{0.12} & \textbf{Yes} \\
\midrule
PS & SP-CF-LA & 0.537 & 0.61 & No \\
PS & SP-CF-AILTA & \textbf{0.041} & \textbf{0.16} & \textbf{Yes} \\
PS & MP-CF & 0.449 & 0.695 & No \\
\bottomrule
\end{tabular}
\end{table}

\subsection{Answer to RQ1}

SP-CF-AILTA significantly outperforms the SP-SR-LA baseline in 4 of 6 domains
(SAT, FlowShop, VRP, PersonnelScheduling; all $p < 0.05$ after Holm correction,
all $A_{12} < 0.20$). In BinPacking it is significantly worse. MP-CF significantly
outperforms the baseline in VRP ($A_{12} = 0.12$) but is significantly worse in SAT.
In BinPacking and TSP, MP-CF's median advantage is not significant.
SP-CF-LA shows no significant improvement over the baseline in any domain and is
significantly worse in BinPacking and VRP.

\textbf{The Choice Function with AILTA acceptance (SP-CF-AILTA) is the only technique
that consistently improves over Simple Random across most domains.}

\section{Results: Cross-Domain Ranking (RQ2)}

\subsection{Rank Table}

Techniques are ranked by median within each domain (1 = best, i.e.\ lowest objective).

\begin{table}[H]
\centering
\begin{tabular}{lcccc}
\toprule
Domain & SP-SR-LA & SP-CF-LA & SP-CF-AILTA & MP-CF \\
\midrule
SAT & 3 & 2 & \textbf{1} & 4 \\
BinPacking & 1.5$^\dagger$ & 3 & 4 & 1.5$^\dagger$ \\
FlowShop & 2 & 3 & \textbf{1} & 4 \\
TSP & 2 & 4 & 3 & \textbf{1} \\
VRP & 3 & 4 & 2 & \textbf{1} \\
PS & 2 & 3 & \textbf{1} & 4 \\
\midrule
\textbf{Mean rank} & \textbf{2.25} & \textbf{3.17} & \textbf{2.00} & \textbf{2.58} \\
\bottomrule
\end{tabular}
\end{table}

$^\dagger$Tied at median 0.0114 (mid-rank); difference not significant ($p = 0.557$).

\subsection{Normalised Scores}

Normalised to $[0, 1]$ per domain: (median $-$ best) / (worst $-$ best). Lower is better.

\begin{table}[H]
\centering
\begin{tabular}{lr}
\toprule
Technique & Mean normalised score \\
\midrule
SP-CF-AILTA & \textbf{0.30} \\
SP-SR-LA & 0.49 \\
MP-CF & 0.50 \\
SP-CF-LA & 0.67 \\
\bottomrule
\end{tabular}
\end{table}

\subsection{Friedman Test}

Friedman chi-squared $= 2.80$, $p = 0.424$. \textbf{Not significant.} There is no
statistically significant difference between the four techniques across all six
domains when treated as a group. This is expected with only 6 blocks (domains) and
relatively close mean ranks (2.00 to 3.17).

\subsection{Answer to RQ2}

SP-CF-AILTA has the best mean rank (2.00) and best mean normalised score (0.30),
ranking 1st in 3 of 6 domains (SAT, FlowShop, PS) and 2nd in VRP. MP-CF ranks
1st in TSP and VRP and ties for 1st in BinPacking. However, the Friedman test
does not detect a significant overall difference ($p = 0.424$), so the ranking
advantage of SP-CF-AILTA over the alternatives cannot be confirmed at the
cross-domain level with 6 domains.

\textbf{No technique dominates across all domains.} SP-CF-AILTA and MP-CF each have
complementary strengths, but MP-CF's advantage is significant only in VRP.

\section{SHAP Analysis (RQ3 and RQ4)}

The XAI analysis addresses the benchmark paper's question---\emph{which heuristics
are effective, and when?}---indirectly: surrogates model acceptance (whether a
heuristic application's result was kept) rather than performance directly. The
features that the surrogate relies on to predict acceptance reveal which aspects of
search state the hyper-heuristic implicitly conditions on when deciding to keep or
discard a candidate solution.

\subsection{Surrogate Fidelity}

Per-domain, per-technique, per-seed gradient boosting surrogates were trained.
\texttt{feat\_llh\_count} was excluded because it is constant within a domain and
acted as a domain label in Assessment~4's pooled surrogates.

\begin{table}[H]
\centering
\begin{tabular}{lrrrl}
\toprule
Technique & Mean ROC-AUC & Min ROC-AUC & Flagged ($<0.6$) & Notes \\
\midrule
SP-SR-LA & 0.52 & 0.49 & \textbf{60/60 (100\%)} & Near-random \\
SP-CF-LA & 0.88 & 0.66 & 0/60 & Adequate \\
SP-CF-AILTA & 0.91 & 0.70 & 0/60 & Good \\
MP-CF & 0.89 & 0.65 & 0/60 & Adequate \\
\bottomrule
\end{tabular}
\end{table}

\textbf{SP-SR-LA surrogates are uniformly flagged across all 6 domains and all 10
seeds} (ROC-AUC 0.49--0.55, mean 0.52). No SHAP attributions are reported for
SP-SR-LA.

\textbf{Robustness check.} Because a uniform AUC of ${\sim}0.5$ could indicate a
logging fault rather than a genuine absence of structure, we verified three potential
failure modes: (1)~\texttt{feat\_last\_h*} columns are populated with real heuristic
indices for SR runs---the distribution across indices is approximately uniform, as
expected from random selection; (2)~acceptance labels are not near-constant---rejection
rates range from 26\% to 40\%, comparable to CF-based techniques;
(3)~\texttt{feat\_iters\_since\_accept} is recorded before the acceptance decision
(the state snapshot is taken before the heuristic is applied and the decision made),
so there is no target leakage.

The trace schema reveals the root cause: \texttt{feat\_last\_h1} records the
heuristic applied in the \emph{previous} iteration, not the current one (the
current heuristic is in the \texttt{heuristic\_id} column, which is not a
surrogate feature). We verified this by checking that
\texttt{feat\_last\_h1} at row~$N$ equals \texttt{heuristic\_id} at
row~$N{-}1$ (100\% match across all traces).

Acceptance rates grouped by the \emph{actual} heuristic applied this iteration
(\texttt{heuristic\_id}) are \emph{not} flat. In BinPacking, local search
heuristics are accepted at ${\sim}100\%$ (they never worsen the solution,
all deltas $\leq 0$). Crossover is also accepted at 100\%, but for a
different reason: in single-point search the incumbent is crossed with
itself, producing an identical child (delta $= 0$ in every application).
One mutation heuristic (h3) also never worsens the solution and is
accepted at 100\%, while the other two mutation heuristics are accepted
at only 0.4--1.9\%. The per-type rates are: local search~100\%,
crossover~100\%, ruin-recreate~57\%, mutation~34\% (the high type
average reflects h3's 100\% rate pulling up the two low ones).

Under random selection, \texttt{feat\_last\_h1} is statistically independent
of the current heuristic, so it carries no signal about the current
acceptance outcome and the surrogate cannot learn from it. Under the
Choice Function, consecutive choices are correlated: CF never repeats
the same heuristic (repeat rate 0\% vs 12.5\% expected under uniform
random), and the $f_2$ pair-sequence term rewards certain orderings.
This temporal structure makes \texttt{feat\_last\_h1} informative about
which heuristic is applied \emph{now}, giving the surrogate a learnable
signal.

We therefore state the finding as: \emph{the SR/CF explainability gap
arises because the lagged heuristic features are uninformative under
random selection but informative under adaptive selection's sequential
structure}.

\subsection{Feature Importance (Ranks and Shares)}

SHAP importance is reported as \textbf{shares} (mean $|\mathrm{SHAP}|$ / total mean
$|\mathrm{SHAP}|$) and \textbf{ranks} (1 = most important). These are comparable
across techniques within a domain. Raw SHAP magnitudes are never compared across
techniques.

\subsubsection*{SP-CF-LA --- Heuristic identity features dominate}

\begin{table}[H]
\centering
\small
\begin{tabular}{llll}
\toprule
Domain & Rank 1 (share) & Rank 2 (share) & Rank 3 (share) \\
\midrule
SAT & \texttt{mean\_logret\_50} (33\%) & \texttt{last\_h1} (21\%) & \texttt{last\_h3} (16\%) \\
BinPacking & \texttt{last\_h3} (27\%) & \texttt{last\_h2} (20\%) & \texttt{last\_class1} (15\%) \\
FlowShop & \texttt{last\_h1} (27\%) & \texttt{last\_h2} (15\%) & \texttt{last\_h3} (13\%) \\
TSP & \texttt{iters\_since\_imp} (32\%) & \texttt{progress} (31\%) & \texttt{mean\_logret\_50} (9\%) \\
VRP & \texttt{last\_h1} (15\%) & \texttt{last\_h2} (13\%) & \texttt{last\_h3} (12\%) \\
PS & \texttt{last\_h1} (18\%) & \texttt{last\_h3} (18\%) & \texttt{iters\_since\_imp} (14\%) \\
\bottomrule
\end{tabular}
\caption*{Feature names are abbreviated: \texttt{last\_h$k$} = \texttt{feat\_last\_h$k$}
(index of the $k$-th most recent heuristic).}
\end{table}

In 4 of 6 domains, the top features are heuristic identity indicators
(\texttt{feat\_last\_h*}). Because these features record the \emph{previous}
iteration's heuristic (\S5.1), they are informative under CF but not under
SR\@. The Choice Function's sequential structure---no consecutive
repeats and $f_2$ pair-sequence rewards---creates temporal correlations
that let the surrogate infer the current heuristic (and hence the
likely acceptance outcome) from the previous one.

\subsubsection*{SP-CF-AILTA --- Search dynamics features dominate}

\begin{table}[H]
\centering
\small
\begin{tabular}{llll}
\toprule
Domain & Rank 1 (share) & Rank 2 (share) & Rank 3 (share) \\
\midrule
SAT & \texttt{progress} (33\%) & \texttt{iters\_since\_imp} (31\%) & \texttt{improve\_rate\_50} (10\%) \\
BinPacking & \texttt{last\_h2} (27\%) & \texttt{last\_h3} (20\%) & \texttt{iters\_since\_acc} (9\%) \\
FlowShop & \texttt{improve\_rate\_50} (50\%) & \texttt{progress} (20\%) & \texttt{mean\_logret\_50} (7\%) \\
TSP & \texttt{iters\_since\_imp} (29\%) & \texttt{progress} (29\%) & \texttt{improve\_rate\_50} (10\%) \\
VRP & \texttt{progress} (30\%) & \texttt{mean\_logret\_50} (20\%) & \texttt{iters\_since\_imp} (16\%) \\
PS & \texttt{last\_h2} (20\%) & \texttt{improve\_rate\_50} (16\%) & \texttt{last\_h1} (14\%) \\
\bottomrule
\end{tabular}
\end{table}

In 4 of 6 domains, the top feature is a search-progress indicator
(\texttt{feat\_progress}, \texttt{feat\_improve\_\allowbreak{}rate\_50},
or \texttt{feat\_iters\_since\_\allowbreak{}improve}).
AILTA's linearly decaying threshold creates a dependency on
\emph{where in the search} the decision occurs, shifting importance away
from heuristic identity toward search dynamics.

\subsubsection*{MP-CF --- Population and temporal features emerge}

\begin{table}[H]
\centering
\small
\begin{tabular}{llll}
\toprule
Domain & Rank 1 (share) & Rank 2 (share) & Rank 3 (share) \\
\midrule
SAT & \texttt{progress} (26\%) & \texttt{iters\_since\_acc} (21\%) & \texttt{pop\_diversity} (18\%) \\
BinPacking & \texttt{progress} (25\%) & \texttt{iters\_since\_imp} (16\%) & \texttt{gap\_to\_best} (10\%) \\
FlowShop & \texttt{improve\_rate\_50} (22\%) & \texttt{mean\_logret\_50} (22\%) & \texttt{pop\_diversity} (16\%) \\
TSP & \texttt{mean\_logret\_50} (36\%) & \texttt{iters\_since\_imp} (17\%) & \texttt{iters\_since\_acc} (17\%) \\
VRP & \texttt{iters\_since\_acc} (26\%) & \texttt{pop\_target\_rank} (22\%) & \texttt{pop\_diversity} (15\%) \\
PS & \texttt{progress} (26\%) & \texttt{iters\_since\_acc} (12\%) & \texttt{iters\_since\_imp} (11\%) \\
\bottomrule
\end{tabular}
\end{table}

Population-specific features (\texttt{feat\_pop\_diversity},
\texttt{feat\_pop\_target\_rank}) appear in the top~3 in 3 of 6 domains
(SAT, FlowShop, VRP). In VRP, where MP-CF ranks 1st in performance,
\texttt{feat\_pop\_target\_rank} has the highest share (22\%)---the surrogate
learned that the rank of the candidate within the population is a strong predictor
of acceptance. The lagged heuristic features (\texttt{feat\_last\_h*}) are consistently
relegated to ranks 9--15 in MP-CF\@. In multi-point search, consecutive
applications typically target different individuals, so the previous
heuristic carries even less signal than in single-point CF\@.

\subsection{Stability Analysis (Kendall's Tau)}

To assess whether SHAP explanations are stable across seeds, we compute Kendall's
$\tau$ rank correlation between all 45 pairs of per-seed feature rankings for each
technique--domain cell.

\begin{table}[H]
\centering
\begin{tabular}{lccc}
\toprule
Domain & SP-CF-LA & SP-CF-AILTA & MP-CF \\
\midrule
SAT & 0.76 (0.08) & 0.59 (0.18) & 0.57 (0.11) \\
BinPacking & 0.38 (0.30) & 0.83 (0.09) & 0.39 (0.14) \\
FlowShop & 0.75 (0.08) & 0.65 (0.10) & 0.42 (0.14) \\
TSP & 0.71 (0.10) & 0.72 (0.10) & 0.42 (0.20) \\
VRP & 0.51 (0.11) & 0.27 (0.17) & 0.68 (0.07) \\
PS & 0.64 (0.14) & 0.72 (0.09) & 0.62 (0.10) \\
\bottomrule
\end{tabular}
\caption*{Values are mean $\tau$ (std). $\tau > 0.6$: moderate--strong agreement; $\tau > 0.8$: strong.}
\end{table}

SP-CF-LA achieves the highest stability in SAT (0.76) and FlowShop (0.75) but is
unstable in BinPacking (0.38). SP-CF-AILTA is most stable in BinPacking (0.83) and
TSP (0.72) but unstable in VRP (0.27). MP-CF is most stable in VRP (0.68) and PS
(0.62) but weaker in BinPacking and FlowShop. No single technique produces
universally stable explanations across all domains.

\subsection{Answer to RQ3}

The features driving acceptance differ by technique and are moderately stable across
seeds ($\tau$ ranges from 0.27 to 0.83 depending on technique--domain cell, with most
cells above 0.5). SP-CF-LA's decisions are explained by lagged heuristic
identity features (acting as a proxy for the current heuristic under CF's
sequential structure), SP-CF-AILTA's by search-progress features, and
MP-CF's by population-level features.
These structural differences between techniques are consistent across domains, even
where per-seed stability within a single domain is lower.

Connecting this to the benchmark question of \emph{which heuristics are effective when}:
the SHAP explanations show that \emph{what the surrogate needs to know} differs by
technique. For SP-CF-LA, heuristic identity is the key signal---CF's sequential
structure makes the previous heuristic predictive of acceptance. For SP-CF-AILTA,
the search stage matters most, indicating that effectiveness is time-dependent. For
MP-CF, population-level context (diversity, rank) matters more than the individual
heuristic applied.

\subsection{Answer to RQ4}

\textbf{Selection strategy}: All 60 SP-SR-LA surrogates have ROC-AUC $\approx 0.52$,
while the 180 CF-based surrogates achieve mean ROC-AUC 0.88--0.91
(min 0.65--0.70). The gap arises because
\texttt{feat\_last\_h*} features record the \emph{previous} iteration's heuristic.
Under random selection these are uninformative noise; under CF's sequential
structure they predict the current heuristic and hence the acceptance outcome
(\S5.1).

\textbf{Acceptance strategy}: Comparing SP-CF-LA (Late Acceptance) with SP-CF-AILTA
(AILTA), the acceptance strategy shifts importance from heuristic identity to search
dynamics. AILTA's linearly decaying threshold makes acceptance sensitive to
\emph{where in the search} the decision is made, causing \texttt{feat\_progress} and
\texttt{feat\_improve\_rate\_50} to dominate.

\textbf{Search topology}: Comparing SP-CF-LA (single-point) with MP-CF (multi-point),
the population topology introduces \texttt{feat\_pop\_diversity} and
\texttt{feat\_pop\_target\_rank} as important features while suppressing heuristic
identity features. In VRP, where MP-CF achieves its strongest significant performance
advantage, \texttt{feat\_pop\_target\_rank} (22\%) and \texttt{feat\_pop\_diversity}
(15\%) together account for 37\% of total SHAP importance.

\section{LIME Analysis}

\subsection{Method}

LIME explanations complement SHAP by providing interpretable \emph{conditions}
(e.g.\ \texttt{feat\_progress <= 0.25}) rather than continuous attributions. For each
surrogate, 50 evenly-spaced instances across \texttt{feat\_progress} are explained.
Each explanation condition is labelled with a sign: positive (+1) means the condition
pushes the surrogate's prediction toward acceptance, negative ($-$1) toward rejection.
Conditions are grouped into 10 progress bands (0--9) to track how acceptance drivers
change over the search.

A condition is classified as an \textbf{acceptance driver} if its
positive\_fraction $> 0.5$ (i.e.\ it more often pushes toward acceptance than
rejection across all explained instances where it appears).

Note that LIME conditions involving heuristic indices (e.g.\
\texttt{3 < feat\_last\_h2 <= 6}) use numeric index ranges. HyFlex exposes each
heuristic's type class via \texttt{getHeuristicsOfType()}. For example, in
BinPacking the mapping is: h0,3,5 = mutation; h1,2 = ruin-recreate;
h4,6 = local search; h7 = crossover. In PersonnelScheduling:
h0--4 = local search; h5--7 = ruin-recreate; h8--10 = crossover;
h11 = mutation. The \texttt{feat\_last\_class*} features,
which record the HyFlex type directly, provide a more interpretable alternative
where they appear in explanations.

\subsection{LIME Acceptance Drivers}

The following summarises the most informative acceptance drivers per technique,
selected from conditions with high positive\_fraction and substantial count.

\subsubsection*{SP-CF-LA --- Heuristic-specific conditions}

\begin{itemize}[nosep]
  \item \texttt{feat\_last\_h3 <= 2.00} (BinPacking: count=144, pos\_frac=0.77, weight=0.098)
  \item \texttt{3.00 < feat\_last\_h2 <= 6.00} (BinPacking: count=132, pos\_frac=0.92, weight=0.165)
  \item \texttt{feat\_last\_h1 > 8.00} (PS: count=127, pos\_frac=1.00, weight=0.112)
  \item \texttt{feat\_progress <= 0.25} (SAT: count=103, pos\_frac=0.95, weight=0.038)
\end{itemize}

SP-CF-LA's acceptance conditions depend on the \emph{previous} iteration's
heuristic. Under CF, the previous heuristic predicts the current one
(\S5.1), matching the SHAP finding that \texttt{feat\_last\_h*} dominates.

\subsubsection*{SP-CF-AILTA --- Search-stage conditions}

\begin{itemize}[nosep]
  \item \texttt{feat\_progress <= 0.25} (SAT: count=130, pos\_frac=1.00, weight=0.263)
  \item \texttt{0.25 < feat\_progress <= 0.50} (SAT: count=120, pos\_frac=1.00, weight=0.240)
  \item \texttt{feat\_iters\_since\_accept\,<=\,0.00}
    (BinPacking: count=364, pos\_frac=1.00, wt=0.129)
  \item \texttt{feat\_improve\_rate\_50 > 0.10} (FlowShop: count=72, pos\_frac=0.90, weight=0.029)
\end{itemize}

The surrogate predicts higher acceptance probability in early progress bands
(progress $\leq 0.50$, positive\_fraction $= 1.0$ across multiple domains).
This is consistent with AILTA's decaying threshold, which is mechanically more
permissive early. However, early-search permissiveness is not unique to AILTA:
SP-CF-LA also shows \texttt{feat\_progress <= 0.25} as an acceptance driver in SAT
(pos\_frac $= 0.95$). Part of the effect may simply reflect that improvements are
easier to find early in any search, regardless of acceptance strategy.

\subsubsection*{MP-CF --- Population diversity conditions}

\begin{itemize}[nosep]
  \item \texttt{feat\_pop\_diversity > 0.05} (SAT: count=53, pos\_frac=0.94, weight=0.308)
  \item \texttt{feat\_pop\_diversity > 0.00} (VRP: count=110, pos\_frac=1.00, weight=0.263)
  \item \texttt{feat\_pop\_diversity > 0.01} (PS: count=70, pos\_frac=0.86, weight=0.282)
  \item \texttt{feat\_gap\_to\_best > 0.04} (BinPacking: count=90, pos\_frac=0.93, weight=0.191)
  \item \texttt{feat\_pop\_target\_rank > 6.00} (VRP: count=26, pos\_frac=1.00, weight=0.202)
\end{itemize}

The surrogate predicts acceptance when population diversity is non-zero
(positive\_fraction 0.86--1.00 with high weights). This is consistent with the
replacement rule: a candidate replaces the worst individual only if strictly better,
and diverse populations offer more weak targets for replacement. The causal direction
is not established---high diversity may be a \emph{consequence} of recent acceptances
rather than a \emph{cause}---but the surrogate reliably uses it as a predictor.

\subsection{LIME Stability}

LIME feature-importance stability is measured by Kendall's $\tau$ on per-seed LIME
share rankings.

\begin{table}[H]
\centering
\begin{tabular}{lccc}
\toprule
Domain & SP-CF-LA & SP-CF-AILTA & MP-CF \\
\midrule
SAT & 0.52 (0.12) & 0.54 (0.13) & 0.47 (0.12) \\
BinPacking & 0.34 (0.23) & 0.69 (0.14) & 0.41 (0.15) \\
FlowShop & 0.39 (0.20) & 0.59 (0.11) & 0.44 (0.12) \\
TSP & 0.40 (0.16) & 0.54 (0.12) & 0.63 (0.09) \\
VRP & 0.21 (0.20) & 0.34 (0.15) & 0.61 (0.11) \\
PS & 0.36 (0.20) & 0.61 (0.10) & 0.49 (0.13) \\
\bottomrule
\end{tabular}
\end{table}

LIME stability is systematically lower than SHAP stability (mean $\tau$ 0.21--0.69
vs 0.27--0.83). This is expected: LIME produces local linear approximations at
individual instances, so it is more sensitive to sampling variation than SHAP's
global TreeExplainer attributions.

\subsection{SHAP vs LIME Agreement}

Both methods identify the same structural patterns:

\begin{enumerate}[nosep]
  \item \textbf{SP-CF-LA}: Heuristic identity features are the primary drivers in both
    SHAP (top shares go to \texttt{feat\_last\_h*}) and LIME (top conditions are
    heuristic-specific ranges).
  \item \textbf{SP-CF-AILTA}: Search-progress features
    (\emph{progress}, \emph{improve\_rate\_50}) dominate both SHAP
    and LIME (progress~$\leq 0.50$, pos.\ frac.\ $=1.0$).
  \item \textbf{MP-CF}: Population features
    (\emph{pop\_diversity}, \emph{pop\_target\_rank}) are key in
    both SHAP and LIME (pos.\ frac.\ ${\approx}\,1.0$).
\end{enumerate}

LIME adds a complementary insight that SHAP cannot provide: the \emph{direction} and
\emph{threshold} at which features drive acceptance predictions. SHAP tells us
\texttt{feat\_progress} is important for SP-CF-AILTA; LIME tells us that
\texttt{feat\_progress <= 0.50} drives acceptance with weight 0.24--0.26.

\section{Limitations}

\begin{enumerate}
  \item \textbf{Not comparable to CHeSC}: This study uses iteration budgets,
    instance~0, and 10 seeds. CHeSC used a normalised 10-minute time limit,
    competition instances, and 31 runs. The real HyFlex domains are the same, but
    the experimental protocol differs.

  \item \textbf{Single instance}: All runs use instance~0. Performance on other
    instances may differ.

  \item \textbf{Friedman test power}: With only 6 domains (blocks), the Friedman test
    has limited power to detect moderate differences in technique rankings.

  \item \textbf{PersonnelScheduling budget}: The 5{,}000-application budget is set by
    convergence analysis but is substantially lower than other domains. This domain's
    results rest on fewer applications.

  \item \textbf{Convergence relative to pilot cap}: BinPacking (94{,}670) and FlowShop
    (91{,}405) reach 95\% convergence very close to the 100{,}000-iteration pilot cap.
    The pilot run may itself not have been long enough, so the true budget needed
    may exceed 100{,}000 and results for these domains may not reflect fully
    converged search.

  \item \textbf{SP-SR-LA unexplainable}: The Simple Random baseline cannot be
    explained by SHAP or LIME because its surrogates fail to learn acceptance
    patterns (ROC-AUC $\approx 0.52$). This is explained by the lagged
    feature encoding: \texttt{feat\_last\_h*} records the previous iteration's
    heuristic, which is uninformative under random selection (\S5.1).
    This limits direct feature-importance comparison to the 3 CF-based
    techniques.

  \item \textbf{LIME stability}: LIME explanations are less stable than SHAP across
    seeds (mean $\tau$ 0.21--0.69 vs 0.27--0.83), particularly for SP-CF-LA in VRP
    ($\tau = 0.21$). Condition-level findings should be interpreted as aggregate
    trends, not per-seed guarantees.

  \item \textbf{Surrogates model acceptance, not performance}: The XAI analysis
    explains which features predict whether a candidate is \emph{accepted}, not
    whether it is \emph{good}. Acceptance is a necessary condition for performance
    improvement but is not sufficient---a technique can accept many poor moves. The
    link between acceptance drivers and technique effectiveness is indirect.
\end{enumerate}

\section{Conclusions}

\begin{enumerate}
  \item \textbf{Replication works}: 240 runs with 10 seeds per cell provide reliable
    variance estimates and enable paired statistical tests.

  \item \textbf{SP-CF-AILTA is the strongest single technique}: It significantly
    outperforms the baseline in 4 of 6 domains, ranks 1st in 3 domains by median,
    and achieves the best mean rank (2.00) and normalised score (0.30).

  \item \textbf{MP-CF shows complementary strengths}: MP-CF achieves the best median
    in TSP and VRP (significant only in VRP) and ties for 1st in BinPacking, suggesting that
    population-based search can help on some problem structures. However, its
    advantage over the baseline is statistically confirmed only in VRP.

  \item \textbf{The SR/CF explainability gap is a feature-encoding effect}:
    All 60 SP-SR-LA surrogates are flagged (ROC-AUC $< 0.6$), while the
    180 CF-based surrogates achieve mean ROC-AUC 0.88--0.91.
    \texttt{feat\_last\_h*} records the \emph{previous} iteration's heuristic.
    Under random selection this is independent of the current move, so
    the surrogate has no signal. Under CF, consecutive choices are
    correlated (no repeats, $f_2$ pair rewards), making the lagged
    feature informative about the current acceptance outcome. Actual
    per-heuristic acceptance rates (grouped by the current
    \texttt{heuristic\_id}) vary widely (e.g.\ 0.4--100\% in BinPacking),
    but the surrogate features do not include the current heuristic
    directly.

  \item \textbf{Each technique type has a distinct explanation signature}: SP-CF-LA is
    explained by lagged heuristic identity, acting as a proxy for the current
    heuristic under CF's sequential structure. SP-CF-AILTA is explained by
    search-progress features, and MP-CF by population-level features. Both
    SHAP and LIME converge on this finding.

  \item \textbf{AILTA's decaying threshold shifts the surrogate's
    reliance to progress features}: The surrogate for SP-CF-AILTA
    consistently ranks \texttt{feat\_progress} and
    \texttt{feat\_improve\_\allowbreak{}rate\_50} above heuristic
    identity features. LIME shows the surrogate predicts higher acceptance
    early in the search, consistent with the linearly decaying threshold.
    However, early-search permissiveness also appears under SP-CF-LA, so
    this pattern may partly reflect general search dynamics rather than
    AILTA specifically.

  \item \textbf{Population diversity predicts acceptance for MP-CF}: The surrogate for
    MP-CF relies on \texttt{feat\_pop\_diversity} and \texttt{feat\_pop\_target\_rank},
    consistent with the replacement rule that targets the worst individual. The causal
    mechanism is not established, but the association is strong in the domains where
    population features reach the top~3 (SAT, FlowShop, VRP).

  \item \textbf{Cross-domain differences are not significant}: The Friedman test
    ($p = 0.424$) does not detect an overall technique effect, likely due to the
    small number of domains.
\end{enumerate}

\end{document}
